\chapter*{Executive Summary}
\addcontentsline{toc}{chapter}{Executive Summary}

\section*{The problem}

Product teams deliberately use separate tools for design, requirements and
implementation, because each is better at its own task than a combined product
would be. The cost is that the relationships between their contents exist only
in the heads of the people doing the work. When a designer moves a screen after
the code was written, or a ticket is closed while the implementation is still
changing, no artifact notices. The failure surfaces in a code review, where a
reviewer must decide whether one pull request still matches the intent it
claims to implement --- a judgement formed across several interfaces, from
memory, with nothing recording what evidence supported it. The reasoning is lost
the moment the change is merged.

\section*{What UNIFY does}

UNIFY subscribes to one design source, one work-item source and one change
source, normalises their events into a single envelope, and stores every
artifact version as an immutable snapshot. Relationships between artifacts are
persisted as typed edges that record how they were established and how confident
the resolver was, so an explicit reference in a pull request title is
distinguishable in the data from an unconfirmed inference. For each atomic
acceptance criterion the system assembles a bounded evidence packet
deterministically, runs a fixed-pass local verification loop over it, and passes
the structured result through a deterministic gate that resolves every citation
against the stored packet. A verdict whose citations do not resolve is
downgraded, and the reason is recorded. The output vocabulary is closed ---
supported, contradicted, partial, or insufficient evidence --- with freshness
and review-required carried as separate fields. All inference runs on
infrastructure the team operates; no repository, ticket or design content is
transmitted to a hosted model interface.

\section*{Research position}

Chapter~2 reviews the comparable systems from their own documentation. CodeRabbit
pulls Jira issue context into a pull request review and returns addressed, not
addressed or unclear per linked objective; Atlassian's Rovo Dev compares a pull
request against the acceptance criteria of a linked Jira work item. Both overlap
directly with part of this project, both are hosted, and neither preserves the
versioned evidence behind a conclusion. The overlap is therefore conceded: the
claimed differentiators are local execution, versioned provenance across tools,
deterministic gating, and reconciliation of late, duplicated and out-of-order
events. Prior research in requirements traceability and automated code review,
and public evaluation corpora, are cited for what they establish rather than
borrowed for authority.

\section*{Status at this evaluation}

\begin{longtable}{p{0.34\textwidth} p{0.22\textwidth} p{0.38\textwidth}}
\caption{Status of the iteration 1 deliverables at the mid evaluation.}\label{tab:exec-status}\\
\hline\textbf{Area} & \textbf{State} & \textbf{Evidence} \\
\hline\endfirsthead
\hline\textbf{Area} & \textbf{State} & \textbf{Evidence} \\
\hline\endhead
Requirements baseline & Complete & 29 functional and 14 non-functional
  requirements, acceptance check per requirement (Ch.~4) \\
Design baseline & Complete & Layered architecture, components, data flow,
  domain model, ERD, verifier loop, anchoring and autonomy-ladder diagrams
  (Ch.~5) \\
Reviewer surface & Complete against typed mock data & 19 routes; typecheck, lint
  and production build pass; 29 static pages generated (Ch.~6) \\
Evaluation-data pilot & Protocol working, annotation incomplete & 15 cases from
  3 TypeScript/React repositories; 138 integrity checks, 0 failures; second
  annotation sheet pending (Ch.~8) \\
Ingestion, evidence graph, rules & Not started & Declared honestly in Ch.~6 and
  Ch.~8; sequenced in Ch.~10 \\
Verifier, gate, local model & Not started & No model downloaded; hardware risk
  still open \\
Process evidence & Maintained & 65 commits, 6 merged pull requests, per-member
  artefacts recorded, GenAI register maintained (Ch.~9) \\
\hline\end{longtable}

\section*{Risks and what comes next}

Three risks carry into the second year. Local inference has never been measured
on this team's hardware, so the ceiling on verification quality is unknown and
the spike is scheduled as early as a stable packet exists. The reliability and
verification work has not started, which is reported as a declared variance
against the iteration plan rather than glossed. And the comparable products have
been read but not trialled, so that comparison rests on documentation until the
dated trials are recorded.

FYP-2 proceeds in dependency order: shared contracts, CI and containers; then
ingestion reliability proven by replay tests; then the evidence graph, analyser
and rules; then the packet builder, bounded verifier and provenance gate, with
the front end bound to the real API and the mock layer removed. Evaluation data
is frozen before any component is tuned against it, two hard gates (rejection of
uncited verdicts, and run-to-run agreement) are enforced before any metric is
reported, and the fine-tuning experiment is decided by a rule written in advance,
so a negative result is reported as such. The frozen scope list in the team
checklist --- no new integrations, no autonomous write-back, no hosted model in
the committed system, no visual certification --- remains in force.
